\section{Finiteness of smooth deformation types}
\label{sec:finiteness}

We now deduce Corollary~\ref{cor:finiteness}. The reduction from an
arbitrary manifold with $b_2\ge5$ to one with a nef isotropic class is
the argument of Sawon~\cite[Proposition~4.3]{Sawon}; we give it here
to show precisely where the second-Betti-number hypothesis enters.
Here deformation
equivalence is generated by smooth proper families of compact
irreducible holomorphic symplectic K\"ahler manifolds over connected
complex bases. The dimension and second Betti number are constant
in such a family.

We use the following smooth consequence of
Engel--Filipazzi--Greer--Mauri--Svaldi~\cite[Theorem~B and
Corollary~A.3]{EFGMS}: in fixed dimension, compact irreducible
holomorphic symplectic K\"ahler manifolds admitting a Lagrangian
fibration have only finitely many smooth deformation types.
Indeed, their Corollary~A.3 deforms such a fibration to a projective
one, and Theorem~B gives finiteness for the projective representatives.
The deformations in these results are locally trivial, meaning
that each point has an analytic neighborhood which deforms as a
product. They therefore preserve smoothness and give ordinary
smooth deformations for the manifolds considered here.

\begin{proof}[Proof of Corollary~\ref{cor:finiteness}]
Let $X$ have complex dimension $2n$ and $r=b_2(X)\ge5$.
Its rational BBF form is nondegenerate and indefinite, of signature
$(3,r-3)$. Meyer's theorem says that an indefinite rational quadratic
form in at least five variables represents zero nontrivially
\cite[Chapter~IV, \S3.2, Corollary~2]{SerreArithmetic}. Clearing denominators and dividing by the positive
divisibility supplies a primitive integral vector
\[
             0\ne e\in H^2(X,\Z),\qquad q(e)=0.
\]

Consider periods perpendicular to $e$ in the marked deformation
component of $X$. The quotient $e^\perp/\R e$ has signature
$(2,r-4)$, so this period locus is nonempty. A generic period in
it is perpendicular to no integral vector outside $\Q e$:
each such additional orthogonality condition is a proper
hyperplane section, and there are countably many of them.
Period surjectivity therefore gives a deformation $X_0$ of $X$
with
\[
                       \NS(X_0)=\Z e.
\]
There are no negative integral $(1,1)$ classes on $X_0$.
The K\"ahler-cone description recalled in
Section~\ref{sec:family} consequently identifies its K\"ahler cone
with its positive cone. One of $e,-e$ lies in the closure of
that cone; the corresponding unique line bundle is nonzero,
nef, and isotropic.

Theorem~\ref{thm:main} makes this bundle semiample and gives a
holomorphic Lagrangian fibration on $X_0$. Thus every smooth
deformation type under consideration contains a manifold with
such a fibration. The smooth consequence of
\cite[Theorem~B and Corollary~A.3]{EFGMS} stated above leaves
only finitely many possibilities in dimension $2n$.
\end{proof}

The restriction $b_2\ge5$ enters only through Meyer's theorem.
Theorem~\ref{thm:main} itself has no such restriction, since its
isotropic class is part of the hypothesis.
